% GitHub Repo and Documentation: https://github.com/celiobjunior/resume-template
% Copyright © 2025 Celio B Junior. All rights reserved.
% 
% Licensed under the Apache License, Version 2.0 (the "License");
% you may not use this file except in compliance with the License.
% You may obtain a copy of the License at
%
%     http://www.apache.org/licenses/LICENSE-2.0
%
% This template follows best practices from README.md

% Início do documento LaTeX: define tipo e formato do documento
% a4paper = tamanho A4, 10pt = tamanho base da fonte
\documentclass[a4paper,10pt]{article}

% --- PACOTES ---
\usepackage[utf8]{inputenc}
\usepackage[T1]{fontenc}
\usepackage[portuguese]{babel}
\usepackage{geometry}
\usepackage{parskip}
\usepackage{hyperref}
\usepackage{titlesec}
% Para quebras de linha em URLs longas (não use \href{}, use \url{} para URLs longas)
% \usepackage{xurl}

% --- CONFIGURAÇÃO DO DOCUMENTO ---
% Define margens da página para maximizar espaço de conteúdo
\geometry{top=1.5cm, bottom=1.5cm, left=1.5cm, right=1.5cm}

% Remove números de página e cabeçalhos para visual limpo do currículo
\pagestyle{empty}

% Metadados do PDF - personalize com suas informações
\hypersetup{
    pdftitle={CV Luis Eduardo Nunes},
    pdfauthor={Luis Eduardo Galvão da Costa Nunes},
    colorlinks=true,
    linkcolor=black,
    urlcolor=black,
    citecolor=black,
    bookmarksdepth=1 
}

% Desabilita numeração das sections
\setcounter{secnumdepth}{0}

% Formata os cabeçalhos de cada section e coloca uma linha em baixo
\titleformat{\section}
{\Large\bfseries}
{}
{0em}
{}
[\titlerule\vspace{0.5ex}]

% --- INÍCIO DO DOCUMENTO ---
\begin{document}

% --- CABEÇALHO ---
\begin{center}
    {\LARGE \textbf{Luis Eduardo Galvão da Costa Nunes}} 
    \\ [0.1cm]
    Natal, RN
    {\textbullet}
    Email: \href{mailto:luisenunes2002@gmail.com}{luisenunes2002@gmail.com} 
    {\textbullet}
    \href{https://www.linkedin.com/in/luisenunes}{linkedin.com/in/luisenunes} 
    {\textbullet}
    \href{https://github.com/LuisENunes}{github.com/LuisENunes}
\end{center}

% --- SEÇÕES ---

\section{Educação}
    \subsection*{\texorpdfstring{
            \textbf{Universidade Potiguar} \hfill Natal, RN
        }{
            Universidade Potiguar -- Natal, RN
        }}
    \textit{Bacharelado em Ciência da Computação \hfill Mar 2025 - 2029 (Cursando)}

% ------

\section{Experiência}
    \subsection*{\texorpdfstring{
            \textbf{Autônomo} \hfill Natal, RN
        }{
            Autônomo -- Natal, RN
        }}
    \textit{Técnico em Manutenção de Computadores \hfill Mar 2023 - Atual}
        \begin{itemize}
            \item Montei centenas de desktops completos para clientes, selecionando componentes e garantindo compatibilidade de hardware
            \item Realizei centenas de formatações e instalações de sistemas operacionais (Windows 10/11), incluindo particionamento e configuração inicial
            \item Diagnostiquei e resolvi centenas de problemas de hardware e software, reduzindo tempo de inatividade para clientes
            \item Executei centenas de limpeza física e remoção de malware, melhorando performance de máquinas lentas
            \item Prestei suporte técnico direto aos clientes, explicando procedimentos e fornecendo orientações de uso
        \end{itemize}

% ------

\section{Certificações}

    \subsection*{\texorpdfstring{
            \textbf{Java Foundations} \hfill Oracle
        }{
            Java Foundations -- Oracle
        }}
        \begin{itemize}
            \item Desenvolvi aplicações utilizando conceitos de Programação Orientada a Objetos (POO), aplicando encapsulamento, herança e polimorfismo para criar código modular e reutilizável
            \item Implementei tratamento de exceções e manipulação de Collections para garantir robustez e eficiência nas aplicações Java
            \item Tech Stack: Java, OOP, POO, Collections, Exception Handling, JVM, Software Development
        \end{itemize}

    \subsection*{\texorpdfstring{
            \textbf{Algoritmos e Lógica de Programação} \hfill Udemy
        }{
            Algoritmos e Lógica de Programação -- Udemy
        }}
        \begin{itemize}
            \item Completei mais de 50 exercícios práticos de lógica de programação, implementando soluções em C, C++, Python, C\# e Java
            \item Apliquei estruturas de dados como vetores e matrizes, utilizando depuração e testes de mesa para validar algoritmos
            \item Tech Stack: Algorithms, Data Structures, Python, C, C++, Problem Solving, Debugging
        \end{itemize}

    \subsection*{\texorpdfstring{
            \textbf{Análise de Dados} \hfill VivaE
        }{
            Análise de Dados -- VivaE
        }}
        \begin{itemize}
            \item Estudei técnicas de coleta, organização e visualização de dados para suporte à tomada de decisão
            \item Apliquei fundamentos de Business Intelligence para interpretação de métricas e geração de insights
            \item Tech Stack: Python, Pandas, NumPy, PowerBI, Data Analysis, Data Visualization, Business Intelligence, Analytics, SQL
        \end{itemize}

% ------

\section{Habilidades}
    \begin{itemize}
        \item \textbf{Técnicas:} Java, Python, JavaScript, HTML, CSS, SQL (MySQL, PostgreSQL, SQL Server), Git, MVC, Montagem e manutenção de hardware
        \item \textbf{Idiomas:} Português (Nativo), Inglês (C1 - Fluente)
    \end{itemize}

\end{document}
